```latex
\documentclass[12pt,a4paper]{article}

% ================= PACOTES =================
\usepackage[utf8]{inputenc}
\usepackage[T1]{fontenc}
\usepackage[brazil]{babel}
\usepackage[margin=2.5cm]{geometry}
\usepackage{graphicx}
\usepackage{booktabs}
\usepackage{array}
\usepackage{amsmath}
\usepackage{amssymb}
\usepackage{hyperref}
\usepackage{longtable}
\usepackage{float}
\usepackage{setspace}
\usepackage{indentfirst}

\hypersetup{
    colorlinks=true,
    linkcolor=black,
    citecolor=black,
    urlcolor=blue,
    pdftitle={PocketBudget - Proposta de Projeto},
    pdfauthor={PocketBudget Team}
}

\setstretch{1.4}

% ================= DOCUMENTO =================
\begin{document}

% ------------------------------------------------------------------
% 1. CAPA
% ------------------------------------------------------------------
\begin{titlepage}
    \centering
    \vspace*{2cm}

    {\large [NOME DA INSTITUIÇÃO]}\par
    \vspace{0.3cm}
    {\large [NOME DO CURSO]}\par

    \vspace{4cm}

    {\Huge \textbf{PocketBudget}}\par
    \vspace{0.5cm}
    {\Large Proposta de Projeto e Documentação Técnica}\par

    \vspace{4cm}

    {\large Álvaro Felipe Farias Pimentel Filho}\par
    {\large Everton Hasabias Furtunato Celestino }\par
    {\large Guilherme Virgílio de Moura Alves}\par
    {\large Matheus Vinícius Almeida Nunes}\par
    {\large Mayck Mitchell Cavalcante da Silva}\par
    {\large Ruan Welisson Nazario Teles}\par
    {\large Vinicius da Silva}\par
    {\large [NOME DO ALUNO 8]}\par

    \vfill

    {\large Professor(a): [NOME DO PROFESSOR]}\par
    \vspace{0.5cm}
    {\large Ano acadêmico: 2026}\par
    \vspace{0.5cm}

    \vspace{1cm}
\end{titlepage}

\tableofcontents
\newpage

% ------------------------------------------------------------------
% 2. INTRODUÇÃO
% ------------------------------------------------------------------
\section{Introdução}

A organização financeira pessoal é um aspecto importante da vida cotidiana, pois influencia diretamente o planejamento, o controle de gastos e a tomada de decisões relacionadas ao uso dos recursos financeiros. Apesar dessa importância, manter um registro consistente de receitas e despesas pode ser uma tarefa pouco prática quando realizada por meio de anotações manuais ou ferramentas genéricas.

O acompanhamento das transações financeiras exige uma forma de registro que seja simples, rápida e acessível. Quando o processo apresenta muitas etapas ou não fornece uma visualização clara das informações registradas, pode ocorrer o abandono do hábito de acompanhamento financeiro.

Nesse contexto, surge o \textbf{PocketBudget}, uma aplicação de gestão financeira pessoal desenvolvida com o objetivo de facilitar o registro e a visualização de receitas e despesas. A aplicação permite cadastrar transações, consultar os valores financeiros registrados, calcular automaticamente o saldo líquido, pesquisar transações e filtrá-las por categoria.

O projeto é composto por duas interfaces: uma aplicação web e uma aplicação mobile nativa para Android. A versão web foi desenvolvida utilizando Angular, enquanto a versão mobile utiliza Kotlin e Jetpack Compose. Ambas possuem mecanismos próprios de persistência local, permitindo que os dados sejam armazenados no dispositivo utilizado pelo usuário sem depender de um servidor externo.

Além de sua finalidade prática, o desenvolvimento do PocketBudget proporciona a aplicação de conceitos relacionados à engenharia de software, desenvolvimento de interfaces, componentização, persistência local, organização de código, controle de versão e documentação técnica.

% ------------------------------------------------------------------
% 3. PROBLEMA
% ------------------------------------------------------------------
\section{Problema}
\label{sec:problema-ref}

O problema central abordado pelo PocketBudget é a dificuldade de registrar e acompanhar transações financeiras pessoais de maneira simples e organizada.

Abordagens manuais, como anotações em papel ou planilhas eletrônicas genéricas, podem apresentar atrito durante o registro de informações. O usuário precisa inserir os dados manualmente e posteriormente organizar ou interpretar essas informações para obter uma visão geral de sua situação financeira.

Aplicativos genéricos de anotações também não possuem necessariamente uma estrutura adequada para operações específicas de controle financeiro, como cálculo automático de saldo, separação entre receitas e despesas, filtragem por categoria e busca por descrição.

O PocketBudget busca reduzir esse atrito por meio de uma interface dedicada ao controle financeiro pessoal. A aplicação concentra as operações essenciais em uma única interface, permitindo registrar transações e visualizar automaticamente os resultados financeiros derivados dos dados cadastrados.

O projeto possui como princípio manter uma solução simples e compatível com o escopo acadêmico proposto, sem introduzir dependências como backend, autenticação ou banco de dados remoto.

% ------------------------------------------------------------------
% 4. OBJETIVOS
% ------------------------------------------------------------------
\section{Objetivos}

\subsection{Objetivo Geral}

Desenvolver uma aplicação de gestão financeira pessoal, denominada PocketBudget, composta por uma interface web e uma interface mobile nativa para Android, permitindo ao usuário registrar receitas e despesas, visualizar seus totais financeiros e localizar transações por meio de busca e filtragem, com persistência local dos dados.

\subsection{Objetivos Específicos}

\begin{itemize}
    \item Desenvolver uma interface web utilizando Angular, TypeScript, HTML e SCSS.
    \item Desenvolver uma aplicação mobile nativa para Android utilizando Kotlin e Jetpack Compose.
    \item Permitir o registro de transações de receita e despesa.
    \item Permitir a exclusão de transações registradas.
    \item Calcular automaticamente o total de receitas.
    \item Calcular automaticamente o total de despesas.
    \item Calcular e exibir o saldo líquido.
    \item Permitir a busca de transações por descrição.
    \item Permitir a filtragem de transações por categoria.
    \item Persistir os dados localmente em cada plataforma.
    \item Utilizar LocalStorage na aplicação web.
    \item Utilizar SharedPreferences com Gson na aplicação mobile.
    \item Aplicar princípios de componentização e separação de responsabilidades.
    \item Utilizar Git e GitHub para controle de versão e colaboração.
    \item Produzir documentação técnica e acadêmica do projeto.
\end{itemize}

% ------------------------------------------------------------------
% 5. METODOLOGIA
% ------------------------------------------------------------------
\section{Metodologia}

O desenvolvimento do PocketBudget foi organizado em etapas relacionadas à definição do problema, levantamento de requisitos, planejamento da interface, definição da arquitetura, implementação e testes.

As principais etapas do desenvolvimento são apresentadas a seguir:

\begin{enumerate}
    \item \textbf{Identificação do problema}: definição da dificuldade relacionada ao registro e acompanhamento de transações financeiras pessoais.
    
    \item \textbf{Definição de requisitos}: identificação das funcionalidades e características necessárias para atender ao objetivo do projeto.
    
    \item \textbf{Ideação}: definição das funcionalidades principais e estabelecimento do escopo da aplicação.
    
    \item \textbf{Wireframing}: elaboração de representações simplificadas das interfaces para auxiliar na definição da organização visual das aplicações.
    
    \item \textbf{Definição da arquitetura}: escolha das tecnologias e organização da estrutura das aplicações web e mobile.
    
    \item \textbf{Implementação}: desenvolvimento das funcionalidades da aplicação web e da aplicação mobile.
    
    \item \textbf{Persistência local}: implementação dos mecanismos de armazenamento necessários para manter as transações após o encerramento das aplicações.
    
    \item \textbf{Testes}: verificação do funcionamento das principais funcionalidades e das regras de negócio.
    
    \item \textbf{Documentação}: organização das informações técnicas e acadêmicas relacionadas ao projeto.
\end{enumerate}

A metodologia prioriza a implementação incremental das funcionalidades, permitindo que diferentes partes do sistema sejam desenvolvidas e verificadas ao longo do processo.

% ------------------------------------------------------------------
% 6. REQUISITOS DO SISTEMA
% ------------------------------------------------------------------
\section{Requisitos do Sistema}

\subsection{Requisitos Funcionais}

\begin{itemize}
    \item RF01 -- O sistema deve permitir o registro de uma transação de receita.
    \item RF02 -- O sistema deve permitir o registro de uma transação de despesa.
    \item RF03 -- O sistema deve permitir a exclusão de uma transação previamente registrada.
    \item RF04 -- O sistema deve permitir a busca de transações por descrição.
    \item RF05 -- O sistema deve permitir a filtragem de transações por categoria.
    \item RF06 -- O sistema deve exibir o total de receitas registradas.
    \item RF07 -- O sistema deve exibir o total de despesas registradas.
    \item RF08 -- O sistema deve exibir o saldo líquido.
    \item RF09 -- O sistema deve persistir os dados das transações localmente.
    \item RF10 -- O sistema deve permitir a visualização da lista de transações registradas.
\end{itemize}

\subsection{Requisitos Não Funcionais}

\begin{itemize}
    \item RNF01 -- \textbf{Responsividade}: a aplicação web deve adaptar sua interface a diferentes tamanhos de tela.
    
    \item RNF02 -- \textbf{Usabilidade}: as interfaces devem apresentar organização simples e intuitiva.
    
    \item RNF03 -- \textbf{Manutenibilidade}: o código deve ser organizado em componentes e estruturas com responsabilidades bem definidas.
    
    \item RNF04 -- \textbf{Desempenho}: as operações de cálculo, busca e filtragem devem apresentar tempo de resposta adequado para o volume de dados esperado.
    
    \item RNF05 -- \textbf{Persistência local}: as aplicações devem manter os dados armazenados localmente após o encerramento e reabertura.
    
    \item RNF06 -- \textbf{Funcionamento offline}: as funcionalidades principais não devem depender de conexão com a internet ou de um servidor remoto.
    
    \item RNF07 -- \textbf{Organização de código}: o projeto deve utilizar uma estrutura de diretórios consistente e facilitar a manutenção do código.
\end{itemize}

% ------------------------------------------------------------------
% 7. ARQUITETURA DO SISTEMA
% ------------------------------------------------------------------
\section{Arquitetura do Sistema}

O PocketBudget é composto por duas aplicações independentes, uma aplicação web e uma aplicação mobile nativa para Android. Ambas implementam funcionalidades equivalentes de controle financeiro, utilizando mecanismos de persistência adequados às suas respectivas plataformas.

\subsection{Aplicação Web}

A aplicação web utiliza Angular como principal framework para construção da interface.

Sua organização é baseada em componentes, serviços, modelos e funções auxiliares. Essa estrutura permite separar responsabilidades e facilitar a manutenção do código.

As principais tecnologias utilizadas são:

\begin{itemize}
    \item \textbf{Angular}: framework utilizado para construção da aplicação web e organização dos componentes da interface.
    
    \item \textbf{TypeScript}: linguagem utilizada na implementação da lógica da aplicação.
    
    \item \textbf{HTML}: linguagem utilizada para estruturação das interfaces dos componentes.
    
    \item \textbf{SCSS}: tecnologia utilizada para estilização das interfaces.
    
    \item \textbf{LocalStorage}: mecanismo de armazenamento disponibilizado pelo navegador para persistência local das transações.
    
    \item \textbf{Componentes Angular}: utilizados para separar funcionalidades como resumo financeiro, formulário e lista de transações.
    
    \item \textbf{Services}: utilizados para centralizar operações relacionadas aos dados das transações.
\end{itemize}

\subsection{Aplicação Mobile}

A aplicação mobile é desenvolvida como uma aplicação nativa para Android utilizando Kotlin e Jetpack Compose.

A interface utiliza o paradigma declarativo do Jetpack Compose, permitindo que os componentes visuais sejam atualizados de acordo com o estado atual da aplicação.

As principais tecnologias utilizadas são:

\begin{itemize}
    \item \textbf{Kotlin}: linguagem utilizada no desenvolvimento da aplicação Android.
    
    \item \textbf{Jetpack Compose}: toolkit utilizado para construção da interface de usuário.
    
    \item \textbf{Material 3}: conjunto de componentes e princípios visuais utilizados na construção da interface.
    
    \item \textbf{ViewModel}: componente utilizado para manter e organizar o estado relacionado aos dados da aplicação.
    
    \item \textbf{SharedPreferences}: mecanismo utilizado para armazenamento local dos dados no dispositivo Android.
    
    \item \textbf{Gson}: biblioteca utilizada para serialização e desserialização dos dados das transações.
\end{itemize}

\subsection{Persistência de Dados}

O PocketBudget utiliza persistência local em ambas as plataformas.

Na aplicação web, os dados são armazenados utilizando o \textit{LocalStorage}, mecanismo disponibilizado pelo navegador para armazenamento persistente de informações no contexto da aplicação.

Na aplicação mobile, os dados são armazenados utilizando \textit{SharedPreferences}. Como as transações são representadas por objetos, a biblioteca Gson é utilizada para converter os dados entre objetos Kotlin e uma representação textual armazenável.

A Tabela~\ref{tab:storage-comparison} apresenta uma comparação entre os mecanismos utilizados.

\begin{table}[H]
    \centering
    \caption{Comparação dos mecanismos de persistência}
    \label{tab:storage-comparison}
    \begin{tabular}{lll}
        \toprule
        \textbf{Característica} & \textbf{Web} & \textbf{Mobile} \\
        \midrule
        Plataforma & Navegador & Android \\
        Tecnologia & LocalStorage & SharedPreferences \\
        Serialização & JSON & Gson \\
        Requer internet & Não & Não \\
        Requer backend & Não & Não \\
        Armazenamento local & Sim & Sim \\
        \bottomrule
    \end{tabular}
\end{table}

% ------------------------------------------------------------------
% 8. MODELO DE DADOS
% ------------------------------------------------------------------
\section{Modelo de Dados}

O modelo de dados central do PocketBudget é representado pela entidade \texttt{Transaction}, responsável por armazenar as informações de cada transação financeira.

A estrutura conceitual é apresentada a seguir:

\begin{verbatim}
Transaction
- id
- description
- amount
- category
- type
- date
\end{verbatim}

Cada campo possui a seguinte finalidade:

\begin{itemize}
    \item \textbf{id}: identificador da transação, utilizado para diferenciar cada registro.
    
    \item \textbf{description}: descrição informada pelo usuário para identificar a transação.
    
    \item \textbf{amount}: valor monetário da transação.
    
    \item \textbf{category}: categoria associada à transação.
    
    \item \textbf{type}: identifica se a transação representa uma receita ou uma despesa.
    
    \item \textbf{date}: data associada ao registro da transação.
\end{itemize}

As categorias utilizadas pelo sistema são:

\begin{itemize}
    \item Food (Alimentação)
    \item Transport (Transporte)
    \item Utilities (Contas)
    \item Entertainment (Entretenimento)
    \item Others (Outros)
\end{itemize}

Os tipos de transação utilizados são:

\begin{itemize}
    \item Income (Receita)
    \item Expense (Despesa)
\end{itemize}

% ------------------------------------------------------------------
% 9. REGRAS DE NEGÓCIO
% ------------------------------------------------------------------
\section{Regras de Negócio}

As principais regras de negócio que orientam o comportamento do PocketBudget são:

\begin{itemize}
    \item RN01 -- O valor de uma transação deve ser positivo.
    
    \item RN02 -- Toda transação deve possuir uma descrição.
    
    \item RN03 -- Toda transação deve possuir uma categoria.
    
    \item RN04 -- Toda transação deve possuir um tipo, podendo ser receita ou despesa.
    
    \item RN05 -- O total de receitas deve considerar somente transações classificadas como receita.
    
    \item RN06 -- O total de despesas deve considerar somente transações classificadas como despesa.
    
    \item RN07 -- O saldo líquido deve ser obtido pela diferença entre o total de receitas e o total de despesas.
    
    \item RN08 -- A exclusão de uma transação deve atualizar os valores apresentados no resumo financeiro.
\end{itemize}

O cálculo do saldo líquido é representado pela Equação~\eqref{eq:saldo-simples}:

\begin{equation}
    B = I - E
    \label{eq:saldo-simples}
\end{equation}

Onde:

\begin{itemize}
    \item $B$ representa o saldo líquido;
    \item $I$ representa o total de receitas;
    \item $E$ representa o total de despesas.
\end{itemize}

Considerando o conjunto de transações registradas, o saldo pode ser representado por:

\begin{equation}
    B = \sum_{i=1}^{n} I_i - \sum_{j=1}^{m} E_j
    \label{eq:saldo-expandido}
\end{equation}

Nessa expressão, $I_i$ representa cada transação de receita e $E_j$ representa cada transação de despesa. Dessa forma, o saldo líquido corresponde à soma das receitas subtraída da soma das despesas.

% ------------------------------------------------------------------
% 10. INTERFACE E EXPERIÊNCIA DO USUÁRIO
% ------------------------------------------------------------------
\section{Interface e Experiência do Usuário}

\subsection{Interface Web}

A aplicação web apresenta uma interface organizada em formato de dashboard, permitindo que as principais informações financeiras sejam visualizadas em uma única página.

Entre os principais elementos da interface estão:

\begin{itemize}
    \item \textbf{Cartões de resumo}: exibem o saldo líquido, o total de receitas e o total de despesas.
    
    \item \textbf{Formulário de transação}: permite cadastrar novas receitas ou despesas.
    
    \item \textbf{Lista de transações}: apresenta as transações registradas pelo usuário.
    
    \item \textbf{Busca}: permite localizar transações utilizando sua descrição.
    
    \item \textbf{Filtro por categoria}: permite restringir a visualização das transações a uma categoria específica.
    
    \item \textbf{Layout responsivo}: adapta a organização dos elementos para diferentes tamanhos de tela.
\end{itemize}

\begin{figure}[H]
    \centering
    % Inserir captura da interface web aqui
    \caption{Interface web do PocketBudget.}
    \label{fig:web-interface}
\end{figure}

\subsection{Interface Mobile}

A aplicação mobile possui as mesmas funcionalidades principais da aplicação web, adaptadas ao contexto de dispositivos Android.

Entre os principais elementos estão:

\begin{itemize}
    \item \textbf{Resumo financeiro}: apresentação dos totais de receitas, despesas e saldo.
    
    \item \textbf{Formulário de transação}: permite cadastrar receitas e despesas por meio de controles adequados à interação por toque.
    
    \item \textbf{Lista de transações}: apresenta os registros armazenados no dispositivo.
    
    \item \textbf{Busca}: permite localizar transações por descrição.
    
    \item \textbf{Filtro por categoria}: permite selecionar categorias específicas para visualização.
    
    \item \textbf{Interface nativa}: utiliza componentes do Jetpack Compose e Material 3 para construção da interface Android.
\end{itemize}

\begin{figure}[H]
    \centering
    % Inserir captura da interface mobile aqui
    \caption{Interface mobile do PocketBudget.}
    \label{fig:mobile-interface}
\end{figure}

% ------------------------------------------------------------------
% 11. ORGANIZAÇÃO DO PROJETO
% ------------------------------------------------------------------
\section{Organização do Projeto}

O repositório do PocketBudget utiliza uma estrutura de monorepo, mantendo as aplicações web, mobile e a documentação acadêmica em um único repositório.

A estrutura principal é organizada da seguinte forma:

\begin{verbatim}
PocketBudget/
├── pocketbudget-web/
├── pocketbudget-mobile/
├── docs/
└── README.md
\end{verbatim}

Os diretórios possuem as seguintes responsabilidades:

\begin{itemize}
    \item \textbf{pocketbudget-web/}: contém o código-fonte da aplicação web desenvolvida com Angular.
    
    \item \textbf{pocketbudget-mobile/}: contém o código-fonte da aplicação mobile nativa desenvolvida com Kotlin e Jetpack Compose.
    
    \item \textbf{docs/}: contém a documentação acadêmica e técnica do projeto.
    
    \item \textbf{README.md}: apresenta informações gerais do projeto, tecnologias utilizadas e instruções para execução das aplicações.
\end{itemize}

A separação das aplicações em diretórios próprios permite que cada plataforma mantenha sua estrutura e ferramentas específicas, enquanto o repositório central facilita o gerenciamento e a colaboração entre os integrantes da equipe.

% ------------------------------------------------------------------
% 12. TECNOLOGIAS
% ------------------------------------------------------------------
\section{Tecnologias}

A Tabela~\ref{tab:tecnologias} apresenta as principais tecnologias utilizadas no desenvolvimento do PocketBudget.

\begin{table}[H]
    \centering
    \caption{Tecnologias utilizadas no desenvolvimento do PocketBudget}
    \label{tab:tecnologias}
    \begin{tabular}{p{3.2cm}p{8.6cm}}
        \toprule
        \textbf{Tecnologia} & \textbf{Papel no projeto} \\
        \midrule
        Angular & Framework utilizado na construção da aplicação web. \\
        
        TypeScript & Linguagem utilizada na implementação da lógica da aplicação web. \\
        
        HTML & Linguagem utilizada para estruturação das interfaces web. \\
        
        SCSS & Tecnologia utilizada para estilização da aplicação web. \\
        
        LocalStorage & Mecanismo utilizado para persistência local das transações na aplicação web. \\
        
        Kotlin & Linguagem utilizada no desenvolvimento da aplicação Android. \\
        
        Jetpack Compose & Toolkit utilizado para construção da interface mobile. \\
        
        Material 3 & Sistema de componentes e padrões visuais utilizado na interface mobile. \\
        
        ViewModel & Componente utilizado para gerenciamento do estado da aplicação mobile. \\
        
        SharedPreferences & Mecanismo utilizado para armazenamento local na aplicação Android. \\
        
        Gson & Biblioteca utilizada para serialização e desserialização dos dados mobile. \\
        
        Git & Sistema de controle de versão utilizado para gerenciamento do código-fonte. \\
        
        GitHub & Plataforma utilizada para hospedagem do repositório e colaboração da equipe. \\
        
        LaTeX & Sistema utilizado para elaboração da documentação acadêmica do projeto. \\
        \bottomrule
    \end{tabular}
\end{table}

% ------------------------------------------------------------------
% 13. CONTROLE DE VERSÃO
% ------------------------------------------------------------------
\section{Controle de Versão}

O controle de versão do PocketBudget é realizado por meio do Git, enquanto o GitHub é utilizado para hospedagem do repositório remoto e colaboração entre os integrantes da equipe.

O uso do Git permite manter um histórico das alterações realizadas no código-fonte e na documentação. Cada alteração pode ser registrada por meio de \textit{commits}, permitindo acompanhar a evolução do projeto durante seu desenvolvimento.

O repositório também permite a utilização de \textit{branches} para desenvolvimento de funcionalidades específicas, possibilitando que alterações sejam realizadas de forma isolada antes de sua integração à branch principal.

O GitHub também facilita a colaboração entre os integrantes da equipe por meio do compartilhamento do código-fonte, gerenciamento de colaboradores e acompanhamento das alterações realizadas no projeto.

% ------------------------------------------------------------------
% 14. TESTES
% ------------------------------------------------------------------
\section{Testes}

Os testes do PocketBudget são direcionados principalmente à verificação das funcionalidades implementadas e das regras de negócio definidas para a aplicação.

As principais verificações realizadas ou previstas incluem:

\begin{itemize}
    \item \textbf{Registro de transações}: verificar se receitas e despesas podem ser cadastradas corretamente.
    
    \item \textbf{Exclusão de transações}: verificar se uma transação pode ser removida e se os valores financeiros são atualizados corretamente.
    
    \item \textbf{Cálculo de totais}: verificar se os valores de receitas, despesas e saldo são calculados corretamente.
    
    \item \textbf{Busca}: verificar se transações podem ser localizadas por meio de sua descrição.
    
    \item \textbf{Filtragem}: verificar se a seleção de uma categoria altera corretamente a lista apresentada.
    
    \item \textbf{Validação}: verificar se dados obrigatórios são exigidos durante o cadastro de uma transação.
    
    \item \textbf{Persistência web}: verificar se os dados continuam disponíveis após o fechamento e reabertura da aplicação web.
    
    \item \textbf{Persistência mobile}: verificar se os dados continuam disponíveis após o fechamento e reabertura da aplicação Android.
    
    \item \textbf{Responsividade}: verificar o comportamento da interface web em diferentes tamanhos de tela.
    
    \item \textbf{Interface mobile}: verificar o comportamento dos componentes da aplicação Android em diferentes tamanhos de tela.
\end{itemize}

A realização dessas verificações busca garantir que as principais funcionalidades do sistema estejam de acordo com os requisitos definidos para o projeto.

% ------------------------------------------------------------------
% 15. CONCLUSÃO
% ------------------------------------------------------------------
\section{Conclusão}

O presente documento apresentou a proposta e a documentação técnica do PocketBudget, uma aplicação de gestão financeira pessoal composta por uma aplicação web e uma aplicação mobile nativa para Android.

O sistema foi desenvolvido com foco na simplicidade de utilização e na disponibilização das funcionalidades essenciais para o controle de receitas e despesas. Entre essas funcionalidades estão o registro e exclusão de transações, cálculo automático de totais, busca por descrição, filtragem por categoria e persistência local dos dados.

Do ponto de vista técnico, a aplicação web utiliza Angular, TypeScript, HTML, SCSS e LocalStorage, enquanto a aplicação mobile utiliza Kotlin, Jetpack Compose, Material 3, ViewModel, SharedPreferences e Gson. A utilização de mecanismos de persistência local permite que as funcionalidades principais sejam executadas sem dependência de backend, autenticação ou banco de dados remoto.

O desenvolvimento das duas aplicações também possibilitou a aplicação prática de conceitos relacionados à componentização, organização de código, gerenciamento de estado, persistência de dados, desenvolvimento responsivo, desenvolvimento mobile nativo, controle de versão e documentação técnica.

A organização do projeto em um repositório único facilita a colaboração entre os integrantes da equipe e mantém as diferentes partes da aplicação separadas de acordo com suas respectivas plataformas.

Por fim, o PocketBudget atende ao escopo definido para o projeto acadêmico ao apresentar uma solução funcional para o registro e acompanhamento de transações financeiras, utilizando tecnologias modernas de desenvolvimento web e mobile e mantendo uma arquitetura compatível com os objetivos e limitações estabelecidos para o trabalho.

\end{document}
```
